% from ega4-21c-fr.tex
\oldpage[IV-4]{319}
Pour avoir un exemple explicite d'anneau parafactoriel de dimension $2$ obtenu de
cette manière et \emph{non factoriel}, considérons l'anneau
\[
  E=\mathbb R[[U,V]]/(U^2+V^2),
\]
dont la clôture intégrale $E'$ s'identifie à $\mathbb C[[U]]$ (6.15.11). Si l'on pose
$A=E[[T]]$, la clôture intégrale de $A$ est l'anneau $A'=E'[[T]]$ (Bourbaki,
\emph{Alg. comm.}, chap.~V, \S1, n\textsuperscript{o}~4, prop.~14); on vérifie
aussitôt que le conducteur de $E$ dans $E'$ est l'idéal maximal $\mathfrak n$ de
$E$, donc le conducteur $\mathfrak J$ de $A$ dans $A'$ est $\mathfrak n[[T]]$, et on
a $B=A/\mathfrak J\cong\mathbb R[[T]]$, $B'=A'/\mathfrak J\cong\mathbb C[[T]]$.
Il est alors immédiat que les conditions $1^{\mathrm o}$, $2^{\mathrm o}$ et
$3^{\mathrm o}$ énoncées ci-dessus sont bien vérifiées, mais $A$ n'est même pas
intégralement clos.

On peut varier cet exemple, et le lecteur verra sans peine que si $k$ est un corps
algébriquement clos, l'anneau $A$, localisé de l'anneau
$k[U,V,W]/(U^2-WV^2)$ en l'idéal maximal engendré par les images de $U$, $V$ et $W$,
vérifie aussi les conditions $1^{\mathrm o}$, $2^{\mathrm o}$ et $3^{\mathrm o}$
ci-dessus.

Il se pourrait que ces trois conditions impliquent même que chacun des anneaux
$B'/\mathfrak p_i'$ soit un anneau de valuation discrète, ce qui entraînerait que
$A'$ est même un anneau \emph{régulier}.
\end{env}

\begin{proposition}[21.13.10]
\label{IV.21.13.10}
Soient $X$ un préschéma localement noethérien, $Y$ une partie fermée de $X$. Pour que
le couple $(X,Y)$ soit parafactoriel il faut et il suffit que, pour tout $y\in Y$,
l'anneau local $\mathcal O_{X,y}$ soit parafactoriel.
\end{proposition}

\begin{proof}
Posons $U=X-Y$ et soit $j:U\to X$ l'injection canonique.

En vertu de (21.13.4), dire que l'homomorphisme canonique
$\mathcal O_X\to j_*(\mathcal O_U)$ soit bijectif équivaut à dire que chacun des
anneaux locaux $\mathcal O_{X,y}$ pour $y\in Y$ satisfait à la condition (i bis) de
(21.13.8).

Montrons d'abord que si le couple $(X,Y)$ est parafactoriel, chacun des anneaux
$\mathcal O_{X,y}$ ($y\in Y$) est parafactoriel. Compte tenu de la remarque
précédente, tout revient à voir que, si l'on pose
$T_y=\operatorname{Spec}(\mathcal O_{X,y})$ et $U_y=T_y-\{y\}$, tout
$\mathcal O_{U_y}$-Module inversible $\mathscr L_0$ est isomorphe à
$\mathcal O_{U_y}$. Or, lorsque $V$ parcourt l'ensemble des voisinages ouverts
affines de $y$ dans $X$, il résulte de (8.2.13) et $(\mathbf I, 2.4.2)$ que le
préschéma $U_y$ est limite projective des préschémas induits par $X$ sur les ouverts
$V-(V\cap\overline{\{y\}})$, qui sont quasi-compacts puisque $X$ est localement
noethérien. Comme les $V-(V\cap\overline{\{y\}})$ sont séparés, il résulte alors de
(8.5.2, (ii)) et (8.5.5) qu'il y a un voisinage ouvert affine $V$ de $y$ dans $X$ et
un $\mathcal O_{U_V}$-Module inversible $\mathscr L'_V$ tel que $\mathscr L_0$ soit le
faisceau induit par $\mathscr L'_V$ sur $U_y$. Mais l'hypothèse entraîne, en vertu de
(21.13.6, (i) et (iii)), que le couple $(V,V\cap\overline{\{y\}})$ est parafactoriel;
on en conclut par définition (21.13.1) qu'il existe un $\mathcal O_V$-Module inversible
$\mathscr L_V$ induisant $\mathscr L'_V$ sur $U_V$. Remplaçant au besoin $V$ par un
voisinage ouvert plus petit de $y$, on peut supposer que $\mathscr L_V$ est isomorphe
à $\mathcal O_V$, d'où on conclut que $\mathscr L_0$ est isomorphe à
$\mathcal O_{U_y}$.

Inversement, prouvons que si tous les $\mathcal O_{X,y}$ ($y\in Y$) sont
parafactoriels le couple $(X,Y)$ est parafactoriel. On est évidemment ramené, compte
tenu de la remarque du début, à montrer que pour tout $\mathcal O_U$-Module inversible
$\mathscr L$, $j_*(\mathscr L)$ est un $\mathcal O_X$-Module inversible (21.13.5).
La question étant locale sur $X$, on peut supposer $X$ noethérien. L'ensemble $V$ des
$x\in X$ tels que la restriction de $j_*(\mathscr L)$ à un voisinage ouvert de $x$
dans $X$ soit inversible est évidemment ouvert dans $X$; il s'agit de montrer que
$V=X$. Pour cela, supposons le contraire et posons $Z=X-V\neq\varnothing$. Soit $y$
un point \emph{maximal} de $Z$; puisque $Z\subset Y$ par définition,
$\mathcal O_{X,y}$ est par hypothèse un anneau parafactoriel. Remplaçant au besoin
$X$ par un voisinage ouvert de $y$ dans $X$, on peut supposer que la restriction de
$j_*(\mathscr L)$
